
\chapter{滤过 $L_\infty$ 代数及滤过 $A_\infty$ 代数的对称化}
\label{app:a3}

在第 7.4.3--7.4.11 小节中，我们用滤过 $L_\infty$ 代数、
$L_\infty$ 模与 $L_\infty$ 同态等概念描述了算子
$\mathfrak p,\mathfrak q,\mathfrak r$。本附录回顾滤过 $L_\infty$
代数等对象的基本定义，并说明滤过 $A_\infty$ 代数的对称化。
对称化是从一个（滤过）$A_\infty$ 代数得到一个（滤过）
$L_\infty$ 代数的典范过程。（$A_\infty$ 代数的对称化在文献中
早已确立；例如见 \textbf{[LaMa95]}。）

设 $C$ 是 $\Lambda_{0,\mathrm{nov}}$ 上的自由分次模。（本附录取
$\Lambda_{0,\mathrm{nov}}=\Lambda^{\mathbb Q}_{0,\mathrm{nov}}$。）
令 $C[1]$ 为移位次数所得的自由分次
$\Lambda_{0,\mathrm{nov}}$-模。考虑 $B_k(C[1])$，并让阶为 $k!$ 的
对称群 $\mathfrak S_k$ 按下式作用其上：
\[
 \sigma(x_1\otimes\cdots\otimes x_k)
 =(-1)^{s(\sigma)}
 x_{\sigma(1)}\otimes\cdots\otimes x_{\sigma(k)},
\]
其中 $\sigma\colon\{1,\ldots,k\}\to\{1,\ldots,k\}$ 是置换，而
\begin{equation}
 s(\sigma)
 =\sum_{\substack{i<j\\\sigma(i)>\sigma(j)}}
 \deg' x_i\,\deg' x_j.
 \tag{A3.1}\label{eq:a3-1}
\end{equation}
这里 $\deg'$ 是移位次数。令
$E_k(C[1])\subset B_k(C[1])$ 为在 $\mathfrak S_k$-作用下不变的
所有元素所成集合。\liftIndex{I139}%
$C[1]$ 的能量滤过在 $E_k(C[1])$ 上诱导能量滤过 $F^\lambda$。
令 $\widehat E(C[1])$ 为
$\bigoplus_{k=0}^{\infty}E_k(C[1])$ 的完备化。

置
\[
 [x_1,\ldots,x_k]
 =\frac1{k!}\sum_{\sigma\in\mathfrak S_k}
 (-1)^{s(\sigma)}
 x_{\sigma(1)}\otimes\cdots\otimes x_{\sigma(k)}.
\]
这些元素生成 $E_k(C[1])$。$B(C[1])$ 上的余代数结构诱导形式余代数
结构
\[
 \Delta\colon\widehat E(C[1])
 \longrightarrow
 \widehat E(C[1])\widehat\otimes\widehat E(C[1]).
\]
注意，$\Delta$ 是分次余交换的。

考虑下列次数为 $+1$ 的运算及其滤过条件：
\[
 \mathfrak l_k\colon E_k(C[1])\longrightarrow C[1],
 \qquad k=0,1,\ldots,
\]
\begin{align}
 \mathfrak l_k\bigl(F^\lambda E_k(C[1])\bigr)
 &\subset F^\lambda(C[1]),
 \tag{A3.2.1}\label{eq:a3-2-1}\\
 \mathfrak l_0(1)
 &\in\bigcup_{\lambda>0}F^\lambda(C[1]).
 \tag{A3.2.2}\label{eq:a3-2-2}
\end{align}
与式~\eqref{eq:3-2-15} 相同，
$\mathfrak l_k\colon E_k(C[1])\to C[1]$ 诱导余导子
\[
 \widehat d_k\colon
 \widehat E(C[1])\longrightarrow\widehat E(C[1]),
\]
只是要用括号 $[x_1,\ldots,x_n]$ 代替张量积。更准确地说，
\[
\begin{aligned}
 \widehat d_k([x_1,\ldots,x_n])
 =\sum_{\operatorname{Sh}(n-k,k)}
 (-1)^*
 \frac{k!(n-k)!}{n!}
 \bigl[
 x_{a_1},\ldots,x_{a_{n-k}},
 \mathfrak l_k([x_{b_1},\ldots,x_{b_k}])
 \bigr].
\end{aligned}
\]
这里求和取遍所有 $(n-k,k)$-洗牌，也就是满足
\[
 \{a_1,\ldots,a_{n-k}\}\cup\{b_1,\ldots,b_k\}
 =\{1,\ldots,n\},\qquad
 a_i<a_{i+1},\quad b_i<b_{i+1}
\]
的 $n$ 元置换，而
\[
 *=\sum_{i=1}^{n-k}\deg' x_{a_i}
 +\sum_{a_i>b_j}\deg' x_{a_i}\deg' x_{b_j}.
\]
\translatornote{源书把这里的 $(n-k,k)$-洗牌称为“长度为 $k$ 的置换”；
但两块合起来排列的是 $\{1,\ldots,n\}$，故应为 $n$ 元置换。见疑误
A-0721。}
仿照引理~\ref{lem:3-2-17}，可以证明
$\widehat d=\sum_k\widehat d_k$ 收敛，并诱导
\[
 \widehat d\colon
 \widehat E(C[1])\longrightarrow\widehat E(C[1]).
\]

\begin{definition}{A3.3}\label{def:a3-3}
\liftIndex{I295}%
若 $\widehat d\circ\widehat d=0$，就称
$\mathfrak l=\{\mathfrak l_k\}_{k\geq0}$ 在 $C$ 上定义了一个
\emph{滤过 $L_\infty$ 代数}结构。若 $\mathfrak l_0=0$，则称该
滤过 $L_\infty$ 代数是\emph{严格的}。
\end{definition}

可以仿照滤过 $A_\infty$ 代数的幺元定义滤过 $L_\infty$ 代数的中心。
具体地，$e$ 属于中心，当且仅当对所有 $k$ 及
$x_1,\ldots,x_{k-1}$ 都有
\[
 \mathfrak l_k([e,x_1,\ldots,x_{k-1}])=0.
\]

还可以仿照定义~\ref{def:3-2-29} 定义滤过 $L_\infty$ 代数为
$G$-有隙的。下文固定同构
\[
 C\cong\overline C\otimes_{\mathbb Q}
 \Lambda^{\mathbb Q}_{0,\mathrm{nov}},
\]
并以
\[
 \overline{\mathfrak l}_k\colon
 E_k(\overline C[1])\longrightarrow\overline C[1]
\]
表示由 $\mathfrak l_k$ 诱导的运算。把系数环限制为 $\mathbb Q$，
便得到一个无滤过 $L_\infty$ 代数。（对无滤过 $L_\infty$ 代数，
我们假设 $\overline{\mathfrak l}_0=0$。）

\begin{example}{A3.4}\label{ex:a3-4}
设 $(\overline C,d,[\,\cdot,\cdot\,])$ 是 $\mathbb Q$ 上的微分分次
李代数。置
\[
 \overline{\mathfrak l}_1(x)=(-1)^{\deg x}dx,\qquad
 \overline{\mathfrak l}_2([x,y])
 =(-1)^{\deg x(\deg y+1)}[x,y],
\]
并令其余所有 $\overline{\mathfrak l}_k=0$。容易验证
$(\overline C,\overline{\mathfrak l}_k)$ 是 $L_\infty$ 代数。
\end{example}

下面定义滤过 $A_\infty$ 代数的对称化。设
$(C,\mathfrak m_*)$ 是滤过 $A_\infty$ 代数。为简单起见，下文以
$\mathfrak l_k(x_1,\ldots,x_k)$ 代替
$\mathfrak l_k([x_1,\ldots,x_k])$。置
\begin{equation}
 \mathfrak l_k(x_1,\ldots,x_k)
 =\frac1{k!}\sum_{\sigma\in\mathfrak S_k}
 (-1)^{s(\sigma)}
 \mathfrak m_k(x_{\sigma(1)},\ldots,x_{\sigma(k)}),
 \tag{A3.5}\label{eq:a3-5}
\end{equation}
其中 $s(\sigma)$ 如式~\eqref{eq:a3-1}，求和取遍
$\{1,\ldots,k\}$ 的置换群中的所有 $\sigma$。显然，
$\mathfrak l_k$ 诱导同态
\[
 \mathfrak l_k\colon E_k(C[1])\longrightarrow C[1].
\]

\begin{proposition}{A3.6}\label{prop:a3-6}
$\mathfrak l_k\colon E_k(C[1])\to C[1]$ 定义一个滤过
$L_\infty$ 代数结构。
\end{proposition}

\begin{proof}
设
\[
 \widehat d\colon\widehat B(C[1])\longrightarrow\widehat B(C[1])
\]
为 $\mathfrak m$ 诱导的余导子。只需证明
\[
 \widehat d(\widehat E(C[1]))\subset\widehat E(C[1]).
\]
事实上，这将蕴含
$\widehat d\colon\widehat E(C[1])\to\widehat E(C[1])$
就是
$\widehat d\colon\widehat B(C[1])\to\widehat B(C[1])$
的限制。

考虑
\[
 \sigma_{ab}\colon\{1,\ldots,\ell\}\longrightarrow
 \{1,\ldots,\ell\},\qquad \sigma_{ab}\in\mathfrak S_\ell,
\]
其中
\[
 \sigma_{ab}(b)=a,\qquad
 \sigma_{ab}(a)=b,\qquad
 \sigma_{ab}(i)=i\quad(i\ne a,b).
\]
令
$\pi_\ell\colon\widehat B(C[1])\to B_\ell(C[1])$ 为投影。计算得
\[
\begin{aligned}
 &(\sigma_{ab}\circ\pi_\ell\circ\widehat d)
 ([x_1,\ldots,x_k])\\
 &\quad=
 \frac1{k!}\sigma_{ab}\Biggl(
 \sum_{\sigma\in\mathfrak S_k}(-1)^{s'(\sigma)}
 \sum_{i=1}^{\ell}
 x_{\sigma(1)}\otimes\cdots\otimes
 \mathfrak m_{k-\ell+1}
 (x_{\sigma(i)},\ldots,x_{\sigma(i+k-\ell)})
 \\
 &\hspace{10em}{}
 \otimes\cdots\otimes x_{\sigma(k)}
 \Biggr).
\end{aligned}
\]
这里
\[
 s'(\sigma)=s(\sigma)
 +\sum_{j=1}^{i-1}\deg' x_{\sigma(j)},
\]
而 $s(\sigma)$ 如式~\eqref{eq:a3-1}。置
\[
\begin{aligned}
 A_i={}&
 \sum_{\sigma\in\mathfrak S_k}(-1)^{s'(\sigma)}
 x_{\sigma(1)}\otimes\cdots\otimes
 \mathfrak m_{k-\ell+1}
 (x_{\sigma(i)},\ldots,x_{\sigma(i+k-\ell)})
 \\
 &\quad{}
 \otimes\cdots\otimes x_{\sigma(k)}.
\end{aligned}
\]
容易看出
\[
 \sigma_{ab}A_a=A_b,\qquad
 \sigma_{ab}A_b=A_a,\qquad
 \sigma_{ab}A_i=A_i\quad(i\ne a,b).
\]
因而
\[
 (\sigma_{ab}\circ\pi_\ell\circ\widehat d)
 ([x_1,\ldots,x_k])
 =(\pi_\ell\circ\widehat d)([x_1,\ldots,x_k]),
\]
这正是所需结论。
\end{proof}

\begin{definition}{A3.7}\label{def:a3-7}
\liftIndex{I096}%
称 $(C,\mathfrak l)$ 为 $(C,\mathfrak m)$ 的\emph{对称化}。无滤过
$A_\infty$ 代数的对称化（它是一个 $L_\infty$ 代数）也可用同样方式
定义。
\end{definition}

\begin{example}{A3.8}\label{ex:a3-8}
设 $(\overline C,d,\wedge)$ 是微分分次代数。置
\begin{equation}
 \overline{\mathfrak m}_1(x)=(-1)^{\deg x}dx,\qquad
 \overline{\mathfrak m}_2(x,y)
 =(-1)^{\deg x(\deg y+1)}(x\wedge y).
 \tag{A3.9}\label{eq:a3-9}
\end{equation}
并令 $\overline{\mathfrak m}_k=0$（$k\ne1,2$），便得到 $A_\infty$ 代数
$(\overline C,\overline{\mathfrak m})$。再以
\[
 [x,y]=\frac12
 \bigl(x\wedge y-(-1)^{\deg x\deg y}y\wedge x\bigr)
\]
定义微分分次李代数
$(\overline C,d,[\,\cdot,\cdot\,])$。容易验证，
$(\overline C,\overline{\mathfrak m})$ 的对称化恰为按例
\ref{ex:a3-4} 从该微分分次李代数得到的 $L_\infty$ 代数。
\end{example}

\begin{proposition}{A3.10}\label{prop:a3-10}
设 $(C,\mathfrak m)$ 是滤过 $A_\infty$ 代数。假设把它的系数约化到
$\mathbb Q$ 所得的
$(\overline C,\overline{\mathfrak m})$，是由微分分次代数
$(\overline C,d,\wedge)$ 按式~\eqref{eq:a3-9} 得到的。再假设
$(\overline C,d,\wedge)$ 分次交换。则对称化
$(C,\mathfrak l)$ 满足
\[
 \mathfrak l_k\equiv0\pmod{\Lambda^+_{0,\mathrm{nov}}},
 \qquad k\ne1.
\]
\end{proposition}

\begin{proof}
只需证明 $\overline{\mathfrak l}_2=0$，因为
$\overline{\mathfrak m}_k=0$（$k\geq3$）显然蕴含
$\overline{\mathfrak l}_k=0$（$k\geq3$）。计算得
\[
\begin{aligned}
 2\overline{\mathfrak l}_2(x,y)
 &=
 \overline{\mathfrak m}_2(x,y)
 +(-1)^{\deg' x\deg' y}\overline{\mathfrak m}_2(y,x)\\
 &=
 (-1)^{\deg x(\deg y+1)}x\wedge y\\
 &\quad{}
 +(-1)^{\deg' x\deg' y+\deg y(\deg x+1)}y\wedge x.
\end{aligned}
\]
由于 $(\overline C,d,\wedge)$ 分次交换，第二项等于
\[
 (-1)^{\deg' x\deg' y+\deg y}x\wedge y
 =(-1)^{\deg x(\deg y+1)+1}x\wedge y,
\]
故与第一项抵消。
\end{proof}

本附录稍后会把命题~\ref{prop:a3-10} 用于
$(C(L;\Lambda_{0,\mathrm{nov}}),\mathfrak m)$ 的典范模型的对称化。

\begin{remark}{A3.11}\label{rem:a3-11}
在本书 2000 年的预印本中，与命题~\ref{prop:a3-10} 证明相同的计算
曾在第 105 页被说成定义滤过 $A_\infty$ 代数对称化时的一个困难。
现在我们明白，这根本不是困难。
\end{remark}

下面定义滤过 $L_\infty$ 同态。设
$(C_i,\mathfrak l^{(i)})$（$i=1,2$）是滤过 $L_\infty$ 代数。
对 $k=0,1,2,\ldots$，考虑下列次数为 $0$ 的映射及其滤过条件：
\[
 \mathfrak f_k\colon E_k(C_1[1])\longrightarrow C_2[1],
\]
\[
 \mathfrak f_k\bigl(F^\lambda E_k(C_1[1])\bigr)
 \subseteq F^\lambda C_2[1].
\]
仿照式~\eqref{eq:3-2-28}，它诱导余代数同态
\[
 \widehat{\mathfrak f}\colon
 \widehat E(C_1[1])\longrightarrow\widehat E(C_2[1]).
\]

\begin{definition}{A3.12}\label{def:a3-12}
\liftIndex{I287}%
若
\[
 \widehat{\mathfrak f}\circ\widehat d^{\,1}
 =\widehat d^{\,2}\circ\widehat{\mathfrak f},
\]
就称 $\mathfrak f=\{\mathfrak f_k\}_{k\geq0}$ 是从 $C_1$ 到 $C_2$
的\emph{滤过 $L_\infty$ 同态}。若 $\mathfrak f_0=0$，则称该同态
是严格的。
\end{definition}

滤过 $L_\infty$ 同态的复合由下式定义：
\[
 \widehat{\mathfrak g\circ\mathfrak f}
 =\widehat{\mathfrak g}\circ\widehat{\mathfrak f}
\]

\begin{lemma}{A3.13}\label{lem:a3-13}
两个滤过 $A_\infty$ 代数之间的任意滤过 $A_\infty$ 同态，经对称化
都会诱导一个滤过 $L_\infty$ 同态。
\end{lemma}

证明很容易，故略去。定义~\ref{def:a3-12} 与引理~\ref{lem:a3-13}
也有无滤过版本和（或）含幺版本。

还可以仿照 $A_\infty$ 情形，为滤过 $L_\infty$ 代数和滤过
$L_\infty$ 同态定义有隙条件。下文假设所有滤过 $L_\infty$ 代数及
滤过 $L_\infty$ 同态都有隙。

下面讨论滤过 $L_\infty$ 同态之间的同伦。设 $(C,\mathfrak l)$ 是
滤过 $L_\infty$ 代数。按定义~\ref{def:4-2-9} 置
\[
 \mathfrak C[1]
 :=\operatorname{Poly}([0,1],C[1])
 \cong\operatorname{Poly}([0,1],C)[1].
\]
再按定义--命题~\ref{def-prop:4-2-15} 的证明，以同样方式定义
\[
 \mathrm{Incl}_k\colon E_k(C[1])\longrightarrow\mathfrak C[1],
 \qquad
 (\mathrm{Eval}_{s=s_0})_k\colon
 E_k(\mathfrak C[1])\longrightarrow C[1]
\]

\begin{lemma}{A3.14}\label{lem:a3-14}
分次 $\Lambda_0$-模
$\mathfrak C=\operatorname{Poly}([0,1],C)$ 上存在一个滤过
$L_\infty$ 代数结构，使 $\mathrm{Incl}$ 与 $\mathrm{Eval}$ 都是滤过
$L_\infty$ 同态。此外，它们在
$\overline{\mathfrak l}_1$-上同调上诱导同构。%
\translatornote{源书把 $\mathrm{Incl}_k$、$\mathrm{Eval}_k$ 的靶分别
写成 $\mathfrak C$、$C$，均漏去滤过 $L_\infty$ 同态分量所需的移位
$[1]$；引理 A3.14 又把承载结构的底模写成
$\operatorname{Poly}([0,1],C[1])$。依定义 A3.12 的型别及定义 4.2.9
的同构
$\operatorname{Poly}([0,1],C)[1]\cong
\operatorname{Poly}([0,1],C[1])$，译文分别补回移位，并以未移位的
$\operatorname{Poly}([0,1],C)$ 为底模。见疑误 A-0723。}
\end{lemma}

证明与引理~\ref{lem:4-2-13} 的证明完全类似，故略去。

\begin{remark}{A3.15}\label{rem:a3-15}
滤过 $A_\infty$ 代数 $C$ 的区间模型还有另一种构造
$C^{[0,1]}$；它似乎不适用于滤过 $L_\infty$ 代数，因为它不保持
对称性（见注~\ref{rem:4-2-21}）。因此，作者不知道怎样在
$\mathbb Z$ 或 $\mathbb Z_2$ 上定义 $L_\infty$ 同态之间的同伦。
$C^{[0,1]}$ 的对称化会成为 $L_\infty$ 代数，但对称化需要在
$\mathbb Q$ 上进行。
\end{remark}

利用模型
$(\operatorname{Poly}([0,1],C),\mathfrak l)$\liftIndex{I362}，可以定义滤过
$L_\infty$ 同态之间的同伦，并证明同伦是与复合相容的等价关系。
于是可以按显然的方式定义滤过 $L_\infty$ 代数的同伦等价。

\begin{theorem}{A3.16}\label{thm:a3-16}\liftIndex{I383}
设 $\mathfrak f=\{\mathfrak f_k\}_{k\geq0}$ 是从 $C_1$ 到 $C_2$
的滤过 $L_\infty$ 同态，并诱导同构
\[
 H(\overline C_1;\overline{\mathfrak l}^{\,1}_1)
 \longrightarrow
 H(\overline C_2;\overline{\mathfrak l}^{\,2}_1).
\]
则存在从 $C_2$ 到 $C_1$ 的滤过 $L_\infty$ 同态 $\mathfrak g$，
使 $\mathfrak f\circ\mathfrak g$ 与
$\mathfrak g\circ\mathfrak f$ 都同伦于恒等同态。
\end{theorem}

本定理的证明与定理~\ref{thm:4-2-45} 的证明平行，故略去。

下面讨论滤过 $L_\infty$ 代数的典范模型。

\begin{definition}{A3.17}\label{def:a3-17}
若 $\overline{\mathfrak l}_1=0$，则称滤过 $L_\infty$ 代数
$(C,\mathfrak l)$ 是\emph{典范的}。
\end{definition}

注意，典范滤过 $A_\infty$ 代数的对称化是典范 $L_\infty$ 代数。

\begin{theorem}{A3.18}\label{thm:a3-18}
任意有隙滤过 $L_\infty$ 代数都同伦等价于一个典范滤过
$L_\infty$ 代数。
\end{theorem}

证明与第~\ref{subsec:5-4-3} 小节中定理~\ref{thm:5-4-2} 的证明
完全类似，故略去。

\begin{theorem}{A3.19}\label{thm:a3-19}
设 $L$ 是 $M$ 的相对 spin 拉格朗日子流形，而
$(H(L;\Lambda_{0,\mathrm{nov}}),\mathfrak m)$ 是推论
\ref{cor:5-4-6} 所得的典范滤过 $A_\infty$ 代数。令
$(H(L;\Lambda_{0,\mathrm{nov}}),\mathfrak l)$ 为它的对称化。
则对所有 $k$ 都有
\begin{equation}
 \mathfrak l_k\equiv0\pmod{\Lambda^+_{0,\mathrm{nov}}}.
 \tag{A3.20}\label{eq:a3-20}
\end{equation}
\end{theorem}

证明需要下述命题。

\begin{proposition}{A3.21}\label{prop:a3-21}
设
\[
 \overline{\mathfrak f}\colon
 (\overline C_1,\overline{\mathfrak m})
 \longrightarrow
 (\overline C_2,\overline{\mathfrak m})
\]
是 $A_\infty$ 代数的同伦等价。设
$(C_1,\mathfrak m)$ 是形变
$(\overline C_1,\overline{\mathfrak m})$ 的有隙滤过
$A_\infty$ 代数。则存在形变
$(\overline C_2,\overline{\mathfrak m})$ 的有隙滤过
$A_\infty$ 代数 $(C_2,\mathfrak m)$，以及满足下列两式的滤过
$A_\infty$ 代数同伦等价 $\mathfrak f$：
\[
\begin{gathered}
 \mathfrak f\colon(C_1,\mathfrak m)
 \longrightarrow(C_2,\mathfrak m),\\
 \mathfrak f\equiv\overline{\mathfrak f}
 \pmod{\Lambda^+_{0,\mathrm{nov}}}.
\end{gathered}
\]
\end{proposition}

命题~\ref{prop:a3-21} 是引理~\ref{lem:7-2-129} 的直接推论。

\begin{proof}[定理~\ref{thm:a3-19} 的证明]
令
\[
 (H^*_{\mathrm{DR}}(L;\mathbb R),
   \overline{\mathfrak l}_{\mathrm{can},\mathbb R})
\]
为德拉姆 D.G.A. $(\Omega(L),d,\wedge)$ 经对称化所得
$L_\infty$ 代数的典范模型。由于楔积是分次交换的，
\[
 \overline{\mathfrak l}_{\mathrm{can},\mathbb R}=0.
\]

另一方面，设
$(H^*(L;\mathbb Q),\overline{\mathfrak l}_{\mathrm{can}})$ 是
$(H^*(L;\Lambda^{\mathbb Q}_{0,\mathrm{nov}}),
\mathfrak l_{\mathrm{can}})$ 的 $\mathbb Q$-约化，则
\[
 (H^*(L;\mathbb Q),\overline{\mathfrak l}_{\mathrm{can}})
 \otimes\mathbb R
\]
同伦等价于
\[
 (H^*_{\mathrm{DR}}(L;\mathbb R),
 \overline{\mathfrak l}_{\mathrm{can},\mathbb R}).
\]
由于两者都是典范的，它们同构。因此
$\overline{\mathfrak l}_{\mathrm{can}}=0$，即得所需结论。
\end{proof}

定理~\ref{thm:a3-19} 蕴含：在拉格朗日子流形的滤过
$A_\infty$ 结构经对称化后，经典杯积与（高阶）Massey 积的贡献
都变为平凡。
